% Chapter 1 – Introduction to the Pravyom Infrastructure Platform
% This file is intended to be included into a larger report or thesis
% that uses the `report` or `book` class.
%
% Example usage in a main document:
%   % Chapter 1 – Introduction to the Pravyom Infrastructure Platform
% This file is intended to be included into a larger report or thesis
% that uses the `report` or `book` class.
%
% Example usage in a main document:
%   % Chapter 1 – Introduction to the Pravyom Infrastructure Platform
% This file is intended to be included into a larger report or thesis
% that uses the `report` or `book` class.
%
% Example usage in a main document:
%   % Chapter 1 – Introduction to the Pravyom Infrastructure Platform
% This file is intended to be included into a larger report or thesis
% that uses the `report` or `book` class.
%
% Example usage in a main document:
%   \include{documentation/pdf/chapter1_pravyom_overview}

\chapter{Introduction to the Pravyom Infrastructure Platform}

\section{Motivation and Context}

Modern distributed systems and blockchain platforms face a persistent tension between theoretical advances and production reality. Many academic systems are evaluated only under narrow benchmarks or with partial implementations; conversely, many industrial systems sacrifice clarity and verifiability in favour of short-term performance. The \emph{Pravyom} platform is designed explicitly to bridge this gap: it is a research-grade, production-validated infrastructure stack that exposes its behaviour through reproducible demos, logs, and proof artefacts.

At its core, Pravyom separates the concerns of community participation and enterprise coordination into two tightly-coupled planes:

\begin{itemize}
  \item The \textbf{BPI plane} (BPI OS, community nodes, and bundles) models how community-operated nodes contribute work, proofs, and resources.
  \item The \textbf{BPCI plane} (BPCI servers, consensus, blockchain, auction infrastructure, and Herm\`es mesh) models how an enterprise infrastructure layer can coordinate, settle, and audit that work at internet scale.
\end{itemize}

This thesis chapter focuses on Pravyom as a coherent system: its architectural principles, its concrete components, and the methodology used to validate that the implementation matches the intended design.

\section{High-Level Architecture}

The Pravyom stack can be viewed as a layered architecture:

\begin{enumerate}
  \item \textbf{Core mathematical and consensus layer}:
    \begin{itemize}
      \item A 6D blockchain and LCCD/QCE2-style consensus framework, designed for strong auditability and quantum-resilient security.
      \item A logbook and PoE (Proof of Evidence) bridge that connects application-level events to immutable 6D headers.
    \end{itemize}
  \item \textbf{BPCI enterprise infrastructure}:
    \begin{itemize}
      \item A constellation of Rust services: consensus server, blockchain server, auction mempool, auction DB maintainer, and a BPI--BPCI bridge.
      \item An autonomous economy module with a four-coin model (GEN/NEX/FLX/AUR) and associated treasury logic.
      \item A web server exposing bank- and government-stamped APIs, registries, and wallet orchestration.
    \end{itemize}
  \item \textbf{Herm\`es P2P mesh and networking}:
    \begin{itemize}
      \item A Hermes-Lite Web-4 mesh providing kappa-aware routing, cellular division, and O(log $n$) discovery behaviour.
      \item A decentralization pathway whereby BPCI servers migrate from a centralized constellation into an autonomous Herm\`es P2P organism.
    \end{itemize}
  \item \textbf{Developer-facing demos and proof artefacts}:
    \begin{itemize}
      \item Short, laptop-safe demos (constellation, lifecycle, mesh, decentralization) implemented as real binaries in the codebase.
      \item Structured reports and logs that capture each demo's behaviour in a machine- and human-readable format.
    \end{itemize}
\end{enumerate}

Importantly, all major components referenced in this work are implemented as real Rust binaries and libraries; there are no mock services in the primary execution paths used for the demonstrations.

\section{Research Questions and Contributions}

This work makes several contributions that are of interest to both systems researchers and advanced practitioners:

\begin{enumerate}
  \item \textbf{End-to-end, reproducible infrastructure design}: We show that it is possible to design a rich blockchain-based infrastructure (with consensus, auction, registry, and P2P mesh subsystems) where \emph{every} demonstration uses production-grade code paths and can be reproduced on commodity hardware.
  \item \textbf{BPI--BPCI separation of concerns}: We formalize a clean separation between the community (BPI) and enterprise (BPCI) planes, including a bridge that carries economic flows, bundles, and proofs, while preserving a transparent lifecycle of tokens and fees.
  \item \textbf{Herm\`es-driven decentralization path}: We present a concrete mechanism by which an initially centralized constellation of BPCI servers can gradually transition into an autonomous Herm\`es P2P mesh, tracked through clearly defined phases (Centralized Constellation, NodeConnectionSync, MeshEvolution, AutonomousMesh).
  \item \textbf{Proof-oriented documentation}: We accompany the implementation with structured proof logs, including command transcripts and result tables, so that a reviewer can not only read about the system but also verify its runtime behaviour step by step.
\end{enumerate}

In combination, these contributions argue that Pravyom is not merely a theoretical architecture but a living, verifiable system that can serve as a reference for next-generation Web3.5 infrastructure.

\section{Methodology and Evidence}

The validation strategy adopted in this work is deliberately conservative. Rather than relying solely on synthetic benchmarks or informal reasoning, we provide:

\begin{itemize}
  \item \textbf{Concrete binaries and commands}: All demos are invoked via explicit \texttt{cargo run} commands, with well-defined environment variables and ports. This ensures that the behaviour described in the text can be reproduced exactly on a reviewer\'s machine.
  \item \textbf{Structured textual reports}: Each demo writes a detailed report to \texttt{/tmp}, summarizing high-level metrics (durations, counts, balances) and low-level traces (per-transaction tables, mesh health snapshots, and phase timelines).
  \item \textbf{Curated proof logs}: Chapter~2 compiles these reports into a single LaTeX- and PDF-format proof log, preserving the raw data while adding explanatory context about what each demo is intended to prove.
\end{itemize}

For a professor or advanced reviewer, this combination of code, commands, and structured logs aims to minimize ambiguity: the system can be assessed both from source and from observed runtime behaviour.

\section{Organization of the Remainder of This Document}

The chapters that follow are organized to separate architectural intent from empirical evidence:

\begin{itemize}
  \item \textbf{Chapter 2 -- BPCI / BPI / Herm\`es Proof Log}: Presents the concrete demos and their logs, including the BPCI constellation 1-minute demo, the BPI\,$\leftrightarrow$\,BPCI lifecycle demo, the Herm\`es mesh demo, and the decentralization phase simulation. Each section describes the commands used, the properties being exercised, and the full textual report generated by the system.
  \item \textbf{Subsequent chapters (not shown here)} may deepen the treatment of formal models (6D blockchain mathematics, consensus guarantees), performance evaluation, or security analysis, depending on the scope of the overall thesis or report.
\end{itemize}

Taken together, Chapter 1 provides the conceptual framing for Pravyom as an infrastructure platform, while Chapter 2 and beyond provide the empirical and formal evidence required for a rigorous academic or professional assessment.


\chapter{Introduction to the Pravyom Infrastructure Platform}

\section{Motivation and Context}

Modern distributed systems and blockchain platforms face a persistent tension between theoretical advances and production reality. Many academic systems are evaluated only under narrow benchmarks or with partial implementations; conversely, many industrial systems sacrifice clarity and verifiability in favour of short-term performance. The \emph{Pravyom} platform is designed explicitly to bridge this gap: it is a research-grade, production-validated infrastructure stack that exposes its behaviour through reproducible demos, logs, and proof artefacts.

At its core, Pravyom separates the concerns of community participation and enterprise coordination into two tightly-coupled planes:

\begin{itemize}
  \item The \textbf{BPI plane} (BPI OS, community nodes, and bundles) models how community-operated nodes contribute work, proofs, and resources.
  \item The \textbf{BPCI plane} (BPCI servers, consensus, blockchain, auction infrastructure, and Herm\`es mesh) models how an enterprise infrastructure layer can coordinate, settle, and audit that work at internet scale.
\end{itemize}

This thesis chapter focuses on Pravyom as a coherent system: its architectural principles, its concrete components, and the methodology used to validate that the implementation matches the intended design.

\section{High-Level Architecture}

The Pravyom stack can be viewed as a layered architecture:

\begin{enumerate}
  \item \textbf{Core mathematical and consensus layer}:
    \begin{itemize}
      \item A 6D blockchain and LCCD/QCE2-style consensus framework, designed for strong auditability and quantum-resilient security.
      \item A logbook and PoE (Proof of Evidence) bridge that connects application-level events to immutable 6D headers.
    \end{itemize}
  \item \textbf{BPCI enterprise infrastructure}:
    \begin{itemize}
      \item A constellation of Rust services: consensus server, blockchain server, auction mempool, auction DB maintainer, and a BPI--BPCI bridge.
      \item An autonomous economy module with a four-coin model (GEN/NEX/FLX/AUR) and associated treasury logic.
      \item A web server exposing bank- and government-stamped APIs, registries, and wallet orchestration.
    \end{itemize}
  \item \textbf{Herm\`es P2P mesh and networking}:
    \begin{itemize}
      \item A Hermes-Lite Web-4 mesh providing kappa-aware routing, cellular division, and O(log $n$) discovery behaviour.
      \item A decentralization pathway whereby BPCI servers migrate from a centralized constellation into an autonomous Herm\`es P2P organism.
    \end{itemize}
  \item \textbf{Developer-facing demos and proof artefacts}:
    \begin{itemize}
      \item Short, laptop-safe demos (constellation, lifecycle, mesh, decentralization) implemented as real binaries in the codebase.
      \item Structured reports and logs that capture each demo's behaviour in a machine- and human-readable format.
    \end{itemize}
\end{enumerate}

Importantly, all major components referenced in this work are implemented as real Rust binaries and libraries; there are no mock services in the primary execution paths used for the demonstrations.

\section{Research Questions and Contributions}

This work makes several contributions that are of interest to both systems researchers and advanced practitioners:

\begin{enumerate}
  \item \textbf{End-to-end, reproducible infrastructure design}: We show that it is possible to design a rich blockchain-based infrastructure (with consensus, auction, registry, and P2P mesh subsystems) where \emph{every} demonstration uses production-grade code paths and can be reproduced on commodity hardware.
  \item \textbf{BPI--BPCI separation of concerns}: We formalize a clean separation between the community (BPI) and enterprise (BPCI) planes, including a bridge that carries economic flows, bundles, and proofs, while preserving a transparent lifecycle of tokens and fees.
  \item \textbf{Herm\`es-driven decentralization path}: We present a concrete mechanism by which an initially centralized constellation of BPCI servers can gradually transition into an autonomous Herm\`es P2P mesh, tracked through clearly defined phases (Centralized Constellation, NodeConnectionSync, MeshEvolution, AutonomousMesh).
  \item \textbf{Proof-oriented documentation}: We accompany the implementation with structured proof logs, including command transcripts and result tables, so that a reviewer can not only read about the system but also verify its runtime behaviour step by step.
\end{enumerate}

In combination, these contributions argue that Pravyom is not merely a theoretical architecture but a living, verifiable system that can serve as a reference for next-generation Web3.5 infrastructure.

\section{Methodology and Evidence}

The validation strategy adopted in this work is deliberately conservative. Rather than relying solely on synthetic benchmarks or informal reasoning, we provide:

\begin{itemize}
  \item \textbf{Concrete binaries and commands}: All demos are invoked via explicit \texttt{cargo run} commands, with well-defined environment variables and ports. This ensures that the behaviour described in the text can be reproduced exactly on a reviewer\'s machine.
  \item \textbf{Structured textual reports}: Each demo writes a detailed report to \texttt{/tmp}, summarizing high-level metrics (durations, counts, balances) and low-level traces (per-transaction tables, mesh health snapshots, and phase timelines).
  \item \textbf{Curated proof logs}: Chapter~2 compiles these reports into a single LaTeX- and PDF-format proof log, preserving the raw data while adding explanatory context about what each demo is intended to prove.
\end{itemize}

For a professor or advanced reviewer, this combination of code, commands, and structured logs aims to minimize ambiguity: the system can be assessed both from source and from observed runtime behaviour.

\section{Organization of the Remainder of This Document}

The chapters that follow are organized to separate architectural intent from empirical evidence:

\begin{itemize}
  \item \textbf{Chapter 2 -- BPCI / BPI / Herm\`es Proof Log}: Presents the concrete demos and their logs, including the BPCI constellation 1-minute demo, the BPI\,$\leftrightarrow$\,BPCI lifecycle demo, the Herm\`es mesh demo, and the decentralization phase simulation. Each section describes the commands used, the properties being exercised, and the full textual report generated by the system.
  \item \textbf{Subsequent chapters (not shown here)} may deepen the treatment of formal models (6D blockchain mathematics, consensus guarantees), performance evaluation, or security analysis, depending on the scope of the overall thesis or report.
\end{itemize}

Taken together, Chapter 1 provides the conceptual framing for Pravyom as an infrastructure platform, while Chapter 2 and beyond provide the empirical and formal evidence required for a rigorous academic or professional assessment.


\chapter{Introduction to the Pravyom Infrastructure Platform}

\section{Motivation and Context}

Modern distributed systems and blockchain platforms face a persistent tension between theoretical advances and production reality. Many academic systems are evaluated only under narrow benchmarks or with partial implementations; conversely, many industrial systems sacrifice clarity and verifiability in favour of short-term performance. The \emph{Pravyom} platform is designed explicitly to bridge this gap: it is a research-grade, production-validated infrastructure stack that exposes its behaviour through reproducible demos, logs, and proof artefacts.

At its core, Pravyom separates the concerns of community participation and enterprise coordination into two tightly-coupled planes:

\begin{itemize}
  \item The \textbf{BPI plane} (BPI OS, community nodes, and bundles) models how community-operated nodes contribute work, proofs, and resources.
  \item The \textbf{BPCI plane} (BPCI servers, consensus, blockchain, auction infrastructure, and Herm\`es mesh) models how an enterprise infrastructure layer can coordinate, settle, and audit that work at internet scale.
\end{itemize}

This thesis chapter focuses on Pravyom as a coherent system: its architectural principles, its concrete components, and the methodology used to validate that the implementation matches the intended design.

\section{High-Level Architecture}

The Pravyom stack can be viewed as a layered architecture:

\begin{enumerate}
  \item \textbf{Core mathematical and consensus layer}:
    \begin{itemize}
      \item A 6D blockchain and LCCD/QCE2-style consensus framework, designed for strong auditability and quantum-resilient security.
      \item A logbook and PoE (Proof of Evidence) bridge that connects application-level events to immutable 6D headers.
    \end{itemize}
  \item \textbf{BPCI enterprise infrastructure}:
    \begin{itemize}
      \item A constellation of Rust services: consensus server, blockchain server, auction mempool, auction DB maintainer, and a BPI--BPCI bridge.
      \item An autonomous economy module with a four-coin model (GEN/NEX/FLX/AUR) and associated treasury logic.
      \item A web server exposing bank- and government-stamped APIs, registries, and wallet orchestration.
    \end{itemize}
  \item \textbf{Herm\`es P2P mesh and networking}:
    \begin{itemize}
      \item A Hermes-Lite Web-4 mesh providing kappa-aware routing, cellular division, and O(log $n$) discovery behaviour.
      \item A decentralization pathway whereby BPCI servers migrate from a centralized constellation into an autonomous Herm\`es P2P organism.
    \end{itemize}
  \item \textbf{Developer-facing demos and proof artefacts}:
    \begin{itemize}
      \item Short, laptop-safe demos (constellation, lifecycle, mesh, decentralization) implemented as real binaries in the codebase.
      \item Structured reports and logs that capture each demo's behaviour in a machine- and human-readable format.
    \end{itemize}
\end{enumerate}

Importantly, all major components referenced in this work are implemented as real Rust binaries and libraries; there are no mock services in the primary execution paths used for the demonstrations.

\section{Research Questions and Contributions}

This work makes several contributions that are of interest to both systems researchers and advanced practitioners:

\begin{enumerate}
  \item \textbf{End-to-end, reproducible infrastructure design}: We show that it is possible to design a rich blockchain-based infrastructure (with consensus, auction, registry, and P2P mesh subsystems) where \emph{every} demonstration uses production-grade code paths and can be reproduced on commodity hardware.
  \item \textbf{BPI--BPCI separation of concerns}: We formalize a clean separation between the community (BPI) and enterprise (BPCI) planes, including a bridge that carries economic flows, bundles, and proofs, while preserving a transparent lifecycle of tokens and fees.
  \item \textbf{Herm\`es-driven decentralization path}: We present a concrete mechanism by which an initially centralized constellation of BPCI servers can gradually transition into an autonomous Herm\`es P2P mesh, tracked through clearly defined phases (Centralized Constellation, NodeConnectionSync, MeshEvolution, AutonomousMesh).
  \item \textbf{Proof-oriented documentation}: We accompany the implementation with structured proof logs, including command transcripts and result tables, so that a reviewer can not only read about the system but also verify its runtime behaviour step by step.
\end{enumerate}

In combination, these contributions argue that Pravyom is not merely a theoretical architecture but a living, verifiable system that can serve as a reference for next-generation Web3.5 infrastructure.

\section{Methodology and Evidence}

The validation strategy adopted in this work is deliberately conservative. Rather than relying solely on synthetic benchmarks or informal reasoning, we provide:

\begin{itemize}
  \item \textbf{Concrete binaries and commands}: All demos are invoked via explicit \texttt{cargo run} commands, with well-defined environment variables and ports. This ensures that the behaviour described in the text can be reproduced exactly on a reviewer\'s machine.
  \item \textbf{Structured textual reports}: Each demo writes a detailed report to \texttt{/tmp}, summarizing high-level metrics (durations, counts, balances) and low-level traces (per-transaction tables, mesh health snapshots, and phase timelines).
  \item \textbf{Curated proof logs}: Chapter~2 compiles these reports into a single LaTeX- and PDF-format proof log, preserving the raw data while adding explanatory context about what each demo is intended to prove.
\end{itemize}

For a professor or advanced reviewer, this combination of code, commands, and structured logs aims to minimize ambiguity: the system can be assessed both from source and from observed runtime behaviour.

\section{Organization of the Remainder of This Document}

The chapters that follow are organized to separate architectural intent from empirical evidence:

\begin{itemize}
  \item \textbf{Chapter 2 -- BPCI / BPI / Herm\`es Proof Log}: Presents the concrete demos and their logs, including the BPCI constellation 1-minute demo, the BPI\,$\leftrightarrow$\,BPCI lifecycle demo, the Herm\`es mesh demo, and the decentralization phase simulation. Each section describes the commands used, the properties being exercised, and the full textual report generated by the system.
  \item \textbf{Subsequent chapters (not shown here)} may deepen the treatment of formal models (6D blockchain mathematics, consensus guarantees), performance evaluation, or security analysis, depending on the scope of the overall thesis or report.
\end{itemize}

Taken together, Chapter 1 provides the conceptual framing for Pravyom as an infrastructure platform, while Chapter 2 and beyond provide the empirical and formal evidence required for a rigorous academic or professional assessment.


\chapter{Introduction to the Pravyom Infrastructure Platform}

\section{Motivation and Context}

Modern distributed systems and blockchain platforms face a persistent tension between theoretical advances and production reality. Many academic systems are evaluated only under narrow benchmarks or with partial implementations; conversely, many industrial systems sacrifice clarity and verifiability in favour of short-term performance. The \emph{Pravyom} platform is designed explicitly to bridge this gap: it is a research-grade, production-validated infrastructure stack that exposes its behaviour through reproducible demos, logs, and proof artefacts.

At its core, Pravyom separates the concerns of community participation and enterprise coordination into two tightly-coupled planes:

\begin{itemize}
  \item The \textbf{BPI plane} (BPI OS, community nodes, and bundles) models how community-operated nodes contribute work, proofs, and resources.
  \item The \textbf{BPCI plane} (BPCI servers, consensus, blockchain, auction infrastructure, and Herm\`es mesh) models how an enterprise infrastructure layer can coordinate, settle, and audit that work at internet scale.
\end{itemize}

This thesis chapter focuses on Pravyom as a coherent system: its architectural principles, its concrete components, and the methodology used to validate that the implementation matches the intended design.

\section{High-Level Architecture}

The Pravyom stack can be viewed as a layered architecture:

\begin{enumerate}
  \item \textbf{Core mathematical and consensus layer}:
    \begin{itemize}
      \item A 6D blockchain and LCCD/QCE2-style consensus framework, designed for strong auditability and quantum-resilient security.
      \item A logbook and PoE (Proof of Evidence) bridge that connects application-level events to immutable 6D headers.
    \end{itemize}
  \item \textbf{BPCI enterprise infrastructure}:
    \begin{itemize}
      \item A constellation of Rust services: consensus server, blockchain server, auction mempool, auction DB maintainer, and a BPI--BPCI bridge.
      \item An autonomous economy module with a four-coin model (GEN/NEX/FLX/AUR) and associated treasury logic.
      \item A web server exposing bank- and government-stamped APIs, registries, and wallet orchestration.
    \end{itemize}
  \item \textbf{Herm\`es P2P mesh and networking}:
    \begin{itemize}
      \item A Hermes-Lite Web-4 mesh providing kappa-aware routing, cellular division, and O(log $n$) discovery behaviour.
      \item A decentralization pathway whereby BPCI servers migrate from a centralized constellation into an autonomous Herm\`es P2P organism.
    \end{itemize}
  \item \textbf{Developer-facing demos and proof artefacts}:
    \begin{itemize}
      \item Short, laptop-safe demos (constellation, lifecycle, mesh, decentralization) implemented as real binaries in the codebase.
      \item Structured reports and logs that capture each demo's behaviour in a machine- and human-readable format.
    \end{itemize}
\end{enumerate}

Importantly, all major components referenced in this work are implemented as real Rust binaries and libraries; there are no mock services in the primary execution paths used for the demonstrations.

\section{Research Questions and Contributions}

This work makes several contributions that are of interest to both systems researchers and advanced practitioners:

\begin{enumerate}
  \item \textbf{End-to-end, reproducible infrastructure design}: We show that it is possible to design a rich blockchain-based infrastructure (with consensus, auction, registry, and P2P mesh subsystems) where \emph{every} demonstration uses production-grade code paths and can be reproduced on commodity hardware.
  \item \textbf{BPI--BPCI separation of concerns}: We formalize a clean separation between the community (BPI) and enterprise (BPCI) planes, including a bridge that carries economic flows, bundles, and proofs, while preserving a transparent lifecycle of tokens and fees.
  \item \textbf{Herm\`es-driven decentralization path}: We present a concrete mechanism by which an initially centralized constellation of BPCI servers can gradually transition into an autonomous Herm\`es P2P mesh, tracked through clearly defined phases (Centralized Constellation, NodeConnectionSync, MeshEvolution, AutonomousMesh).
  \item \textbf{Proof-oriented documentation}: We accompany the implementation with structured proof logs, including command transcripts and result tables, so that a reviewer can not only read about the system but also verify its runtime behaviour step by step.
\end{enumerate}

In combination, these contributions argue that Pravyom is not merely a theoretical architecture but a living, verifiable system that can serve as a reference for next-generation Web3.5 infrastructure.

\section{Methodology and Evidence}

The validation strategy adopted in this work is deliberately conservative. Rather than relying solely on synthetic benchmarks or informal reasoning, we provide:

\begin{itemize}
  \item \textbf{Concrete binaries and commands}: All demos are invoked via explicit \texttt{cargo run} commands, with well-defined environment variables and ports. This ensures that the behaviour described in the text can be reproduced exactly on a reviewer\'s machine.
  \item \textbf{Structured textual reports}: Each demo writes a detailed report to \texttt{/tmp}, summarizing high-level metrics (durations, counts, balances) and low-level traces (per-transaction tables, mesh health snapshots, and phase timelines).
  \item \textbf{Curated proof logs}: Chapter~2 compiles these reports into a single LaTeX- and PDF-format proof log, preserving the raw data while adding explanatory context about what each demo is intended to prove.
\end{itemize}

For a professor or advanced reviewer, this combination of code, commands, and structured logs aims to minimize ambiguity: the system can be assessed both from source and from observed runtime behaviour.

\section{Organization of the Remainder of This Document}

The chapters that follow are organized to separate architectural intent from empirical evidence:

\begin{itemize}
  \item \textbf{Chapter 2 -- BPCI / BPI / Herm\`es Proof Log}: Presents the concrete demos and their logs, including the BPCI constellation 1-minute demo, the BPI\,$\leftrightarrow$\,BPCI lifecycle demo, the Herm\`es mesh demo, and the decentralization phase simulation. Each section describes the commands used, the properties being exercised, and the full textual report generated by the system.
  \item \textbf{Subsequent chapters (not shown here)} may deepen the treatment of formal models (6D blockchain mathematics, consensus guarantees), performance evaluation, or security analysis, depending on the scope of the overall thesis or report.
\end{itemize}

Taken together, Chapter 1 provides the conceptual framing for Pravyom as an infrastructure platform, while Chapter 2 and beyond provide the empirical and formal evidence required for a rigorous academic or professional assessment.
